% Part: methods
% Chapter: induction
% Section: relations

\documentclass[../../../include/open-logic-section]{subfiles}

\begin{document}

\olfileid{mth}{ind}{rel}

\olsection{సంబంధాలు మరియు ప్రమేయాలు}

వస్తువుల సమితిని (సహజ సంఖ్యలు లేదా చక్కని పదాలు వంటివాటిని) ఆగమనాత్మకంగా నిర్వచించిన తర్వాత, ఆ వస్తువుల\emph{పై సంబంధాలను} కూడా ఆగమనం ద్వారా నిర్వచించవచ్చు. ఉదాహరణకు, ఒక చక్కని పదం~$t_1$, మరో చక్కని పదం~$t_2$లో భాగంగా కనిపిస్తే దాన్ని రెండవదాని ఉపపదం అందాం. దీనికి $t_1 \sqsubseteq t_2$ అనే సంకేతాన్ని వాడుదాం. ప్రతి చక్కని పదమూ తనకే ఉపపదం: $t \sqsubseteq t$. ఈ సంబంధాన్ని ఆగమనాత్మకంగా ఇలా నిర్వచించవచ్చు:

\begin{defn}
  చక్కని పదం~$t_1$,~$t_2$ యొక్క ఉపపదం కావడం అనే సంబంధాన్ని, $t_1
  \sqsubseteq t_2$ను,~$t_2$పై ఆగమనం ద్వారా ఇలా నిర్వచిస్తాం:
  \begin{enumerate}
  \item $t_2$ అక్షరమైతే, $t_1 \sqsubseteq t_2$ అప్పుడూ అప్పుడే $t_1 = t_2$.
  \item $t_2$ అనేది $[s_1 \circ s_2]$ అయితే, $t_1 \sqsubseteq t_2$ అప్పుడూ అప్పుడే $t_1 =
    t_2$, లేదా $t_1 \sqsubseteq s_1$, లేదా $t_1 \sqsubseteq s_2$.
  \end{enumerate}
\end{defn}

ఉదాహరణకు, ఈ నిర్వచనం $\mathrm{a}
\sqsubseteq [\mathrm{b} \circ \mathrm{a}]$ అని చెబుతుంది. ఎందుకంటే (2) ప్రకారం $\mathrm{a} \sqsubseteq [\mathrm{b} \circ \mathrm{a}]$ అప్పుడూ అప్పుడే $\mathrm{a}
= [\mathrm{b} \circ \mathrm{a}]$, లేదా $\mathrm{a} \sqsubseteq b$, లేదా $\mathrm{a} \sqsubseteq \mathrm{a}$. మొదటి రెండు అసత్యం: $\mathrm{a}$ అనేది $[\mathrm{b} \circ
  \mathrm{a}]$కు సమానం కాదు; (1) ప్రకారం $\mathrm{a} \sqsubseteq \mathrm{b}$ అప్పుడూ అప్పుడే $\mathrm{a} = \mathrm{b}$, ఇదీ అసత్యమే. అయితే (1) ప్రకారమే $\mathrm{a} \sqsubseteq \mathrm{a}$ అప్పుడూ అప్పుడే $\mathrm{a} = \mathrm{a}$; ఇది సత్యం.

ఈ నిర్వచనం సఫలం కావడం మనం ఇంకా నిరూపించని ఒక విషయంపై ఆధారపడుతుందని గమనించాలి: ప్రతి చక్కని పదం~$t$ ఒక్క అక్షరం కావాలి, లేదా $t = [s_1 \circ s_2]$ అయ్యే విధంగా \emph{ఏకైకంగా నిర్ణయించబడే} చక్కని పదాలు $s_1$, $s_2$ ఉండాలి. ఇక్కడ ``ఏకైకంగా నిర్ణయించబడే'' అంటే, $t = [s_1 \circ s_2]$ అయితే ఆ పదం, $s_1 \neq r_1$ లేదా $s_2 \neq r_2$ అయ్యేలా \emph{మళ్లీ} $= [r_1 \circ r_2]$ కావద్దు. అలా అయితే (2)వ నిబంధన తనతోనే విరోధపడవచ్చు: $t_2$ను $[s_1 \circ s_2]$గా చదివితే $t_1
\sqsubseteq t_2$ రావచ్చు; కానీ $t_2$ను $[r_1 \circ r_2]$గా చదివితే $t_1 \sqsubseteq t_2$ కాదని రావచ్చు. అలా జరగదని నిరూపించే ముందు, అలా \emph{జరగగల} ఒక ఉదాహరణ చూద్దాం.

\begin{defn}
  \emph{బ్రాకెట్లు లేని పదాలను} ఆగమనాత్మకంగా ఇలా నిర్వచించండి:
  \begin{enumerate}
  \item ప్రతి అక్షరం బ్రాకెట్లు లేని పదం.
    \item $s_1$, $s_2$ బ్రాకెట్లు లేని పదాలైతే, $s_1 \circ s_2$ కూడా బ్రాకెట్లు లేని పదం.
    \item మరేదీ బ్రాకెట్లు లేని పదం కాదు.
  \end{enumerate}
\end{defn}

బ్రాకెట్లు లేని పదాలకు $\mathrm{a}$, $\mathrm{b} \circ
\mathrm{d}$, $\mathrm{b} \circ \mathrm{a} \circ \mathrm{b}$ ఉదాహరణలు. వాటికి ``ఉపపదం''ను పైన చెప్పినట్లే నిర్వచిస్తే, రెండవ నిబంధన ఇలా ఉంటుంది:
\begin{center}
  $t_2 = s_1 \circ s_2$ అయితే, $t_1 \sqsubseteq t_2$ అప్పుడూ అప్పుడే $t_1 = t_2$, లేదా $t_1
  \sqsubseteq s_1$, లేదా $t_1 \sqsubseteq s_2$.
\end{center}

ఇప్పుడు $\mathrm{b} \circ \mathrm{a} \circ \mathrm{b}$ అనేది $s_1
\circ s_2$ రూపంలో ఇలా ఉంటుంది:
\begin{align*}
  s_1 & = \mathrm{b} \text{ మరియు} & s_2 & = \mathrm{a} \circ \mathrm{b}.
\intertext{అది $r_1 \circ r_2$ రూపంలో కూడా ఇలా ఉంటుంది:}
  r_1 & = \mathrm{b} \circ \mathrm{a} \text{ మరియు} & r_2 &= \mathrm{b}.
\end{align*}
అయితే $\mathrm{a} \circ \mathrm{b}$ అనేది $\mathrm{b} \circ
\mathrm{a} \circ \mathrm{b}$కు ఉపపదమా? మొదటి విధంగా చదివితే అవును; రెండవ విధంగా చదివితే కాదు.

ఒక చక్కని పదాన్ని ఇతర చక్కని పదాల నుంచి నిర్మించే విధానం ఏకైకంగా ఉండడాన్ని \emph{ఏకైక పఠనీయత} అంటారు. అటువంటి ఆగమనాత్మకంగా నిర్వచించిన వస్తువులపై సంబంధాల ఆగమనాత్మక నిర్వచనాలు ముఖ్యమైనవి కనుక, ఈ లక్షణం ఉందని నిరూపించాలి.

\begin{prop}
  $t$ ఒక చక్కని పదమనుకుందాం. అప్పుడు $t$ ఒక్క అక్షరం, లేదా $t =
  [s_1 \circ s_2]$ అయ్యే విధంగా ఏకైకంగా నిర్ణయించబడే చక్కని పదాలు $s_1$, $s_2$ ఉంటాయి.
\end{prop}

\begin{proof}
  $t$ ఒక్క అక్షరమైతే షరతు నెరవేరుతుంది. కనుక $t$ ఒక్క అక్షరం కాదనుకుందాం. అప్పుడు ఆగమనాత్మక నిర్వచనం ప్రకారం ఏవో చక్కని పదాలు $s_1$,~$s_2$ కోసం $t$ అనేది $[s_1 \circ s_2]$ రూపంలో ఉండాలి. ఇవి ఏకైకంగా నిర్ణయించబడతాయని చూపడం మిగిలింది: $t = [r_1 \circ r_2]$ అయితే $s_1 = r_1$, $s_2 = r_2$ అని చూపాలి.

  అయితే చక్కని పదాలు $s_1$, $s_2$, $r_1$, $r_2$ కోసం $t = [s_1 \circ s_2]$, అలాగే $t = [r_1 \circ r_2]$ అనుకుందాం. $s_1 = r_1$, $s_2 = r_2$ అని చూపాలి. మొదట $s_1$, $r_1$ ఒకటే కావాలి; కాకపోతే ఒకటి మరొకదానికి నిజ ప్రారంభ భాగం అవుతుంది. కానీ $s_1$,~$r_1$ రెండూ చక్కని పదాలైతే, \olref[sti]{prop:initial} ప్రకారం అది అసాధ్యం. ఇక $s_1 = r_1$ అయితే స్పష్టంగా $s_2 = r_2$ కూడా.
\end{proof}

ప్రమేయాలను కూడా ఆగమనాత్మకంగా నిర్వచించవచ్చు. ఉదాహరణకు, ఏ చక్కని పదాన్నైనా అందులోని $[\dots]$ గూళ్ల గరిష్ఠ లోతుకు పంపే ప్రమేయం~$f$ను ఇలా నిర్వచించవచ్చు:

\begin{defn}
  \ollabel{defn:depth} చక్కని పదం యొక్క \emph{లోతు}, $f(t)$ను, ఆగమనాత్మకంగా ఇలా నిర్వచిస్తాం:
  \[
  f(t) = \begin{cases}
    0 & \text{ $t$ అక్షరమైతే}\\
    \max(f(s_1), f(s_2)) + 1 & \text{ $t = [s_1 \circ s_2]$ అయితే.}
  \end{cases}
  \]
\end{defn}

ఉదాహరణకు,
\begin{align*}
  f([\mathrm{a} \circ \mathrm{b}]) & =
    \max(f(\mathrm{a}),f(\mathrm{b})) + 1 = \\ 
  &= \max(0, 0) + 1 = 1, \text{ మరియు}\\
  f([[\mathrm{a} \circ \mathrm{b}] \circ \mathrm{c}]) & =
    \max(f([\mathrm{a} \circ \mathrm{b}]), f(\mathrm{c})) + 1 = \\ 
  & = \max(1,0) + 1 = 2.
\end{align*}

ఇక్కడ $s_1$, $s_2$ చక్కని పదాలని తీసుకొని, ప్రతి చక్కని పదమూ ఒక్క అక్షరం లేదా $[s_1 \circ s_2]$ రూపంలో ఉంటుందనే విషయాన్ని వాడుతున్నాం. ఆ రూపంలో ఉండడానికి ఒకే మార్గం ఉండడం కూడా ముఖ్యం. ఎందుకో చూడడానికి, ముందు నిర్వచించిన బ్రాకెట్లు లేని పదాలను మళ్లీ పరిశీలిద్దాం. వాటికి అనురూపమైన ``నిర్వచనం'' ఇలా ఉంటుంది:
\[
  g(t) =
  \begin{cases}
    0 & \text{ $t$ అక్షరమైతే}\\
   \max(g(s_1), g(s_2)) + 1 & \text{ $t = s_1 \circ s_2$ అయితే.}
  \end{cases}
\]
ఇప్పుడు బ్రాకెట్లు లేని $\mathrm{a} \circ \mathrm{b} \circ
\mathrm{c} \circ \mathrm{d}$ పదాన్ని చూడండి. దీన్ని ఒకటి కంటే ఎక్కువ విధాలుగా చదవవచ్చు; ఉదాహరణకు $s_1 \circ s_2$గా ఇలా:
\begin{align*}
  s_1 & = \mathrm{a} \text{ మరియు} &
  s_2 & = \mathrm{b} \circ \mathrm{c} \circ \mathrm{d},
\intertext{లేదా $r_1 \circ r_2$గా ఇలా:}
  r_1 & = \mathrm{a} \circ b \text{ మరియు} &
  r_2 &= \mathrm{c} \circ \mathrm{d}.
\end{align*}
మొదటి పఠనం ప్రకారం $g$ను గణిస్తే
\begin{align*}
  g(s_1 \circ s_2) & =
  \max(g(\mathrm{a}), g(\mathrm{b} \circ \mathrm{c} \circ \mathrm{d})) + 1 =\\ & =
  \max(0,2) + 1 = 3
  \intertext{మరో పఠనం ప్రకారం అయితే}
  g(r_1 \circ r_2) & =
  \max(g(\mathrm{a} \circ \mathrm{b}), g(\mathrm{c} \circ \mathrm{d})) + 1=\\ &
  = \max(1,1) + 1 = 2
\end{align*}
కానీ ప్రమేయం ఎల్లప్పుడూ ఏకైక విలువనే ఇవ్వాలి; కాబట్టి $g$కు ఇచ్చిన మన ``నిర్వచనం'' అసలు ప్రమేయాన్ని నిర్వచించదు.

\begin{prob}
  $l$ అనే ప్రమేయానికి ఆగమనాత్మక నిర్వచనం ఇవ్వండి; ఇక్కడ $l(t)$ అనేది చక్కని పదం~$t$లోని చిహ్నాల సంఖ్య.
\end{prob}

\begin{prob}
  చక్కని పదాలు $t$పై నిర్మాణాత్మక ఆగమనం ద్వారా $f(t) < l(t)$ అని నిరూపించండి (ఇక్కడ $l(t)$ అనేది~$t$లోని చిహ్నాల సంఖ్య; $f(t)$ అనేది \olref[mth][ind][rel]{defn:depth}లో నిర్వచించిన $t$ లోతు).
\end{prob}

\end{document}
